\documentclass[aspectratio=169,11pt]{beamer}
\usepackage{beamerucad}
\usepackage{listings}
\definecolor{ckw}{HTML}{2563AB}\definecolor{cst}{HTML}{2E8B57}\definecolor{ccm}{HTML}{8A8F98}
\definecolor{outbg}{HTML}{F4F5F7}\definecolor{outrule}{HTML}{2E8B57}
\lstdefinestyle{py}{language=Python,basicstyle=\ttfamily\scriptsize,
 keywordstyle=\color{ckw}\bfseries,stringstyle=\color{cst},commentstyle=\color{ccm}\itshape,
 showstringspaces=false,columns=fullflexible,keepspaces=true,
 backgroundcolor=\color{ucadband!8},frame=leftline,framerule=2.5pt,rulecolor=\color{ucadband},
 xleftmargin=8pt,aboveskip=3pt,belowskip=1pt,basewidth=0.5em}
\lstdefinestyle{out}{basicstyle=\ttfamily\scriptsize\color{black!82},
 backgroundcolor=\color{outbg},frame=leftline,framerule=2.5pt,rulecolor=\color{outrule},
 xleftmargin=8pt,aboveskip=2pt,belowskip=1pt,columns=fullflexible,keepspaces=true}
\lstset{style=py}
\newcommand{\resultat}{{\scriptsize\textbf{R\'esultat}~$\rightarrow$}\par\vspace{1pt}}

\title{Programmation Python — Séance 5}
\subtitle{Programmation orientée objet}
\author{Dr. El Hadji Bassirou TOURÉ}
\institute{Département de Mathématiques et Informatique\\ Faculté des Sciences et Techniques\\ Université Cheikh Anta Diop de Dakar}
\date{}
\setucadfoot{Programmation Python — Séance 5}
\begin{document}

\begin{frame}[plain,noframenumbering]\titlepage\end{frame}

% =====================================================================
\begin{frame}{Objectifs de la séance}
\begin{ucaddef}[Contenu]
{\small Jusqu'ici, les \emph{données} (variables) et les \emph{traitements} (fonctions) étaient séparés. La programmation orientée objet les \textbf{réunit} : un objet regroupe des données et les comportements qui les manipulent. C'est le style dans lequel sont écrits \texttt{numpy}, \texttt{pandas} et la plupart des bibliothèques.}
\end{ucaddef}
\begin{itemize}\small
  \item \textbf{Classe} et \textbf{objet} : \texttt{class}, \texttt{\_\_init\_\_}, \texttt{self}, attributs.
  \item \textbf{Méthodes} : les comportements de l'objet ; \texttt{\_\_str\_\_} pour l'affichage.
  \item \textbf{Héritage} : spécialiser une classe sans la réécrire.
  \item \textbf{Polymorphisme} : un même appel, un comportement propre à chaque classe.
\end{itemize}
{\small L'entraînement se fait dans le \textbf{notebook}, avec une cellule « À vous » après chaque notion.}
\end{frame}

\begin{frame}{Deux mécanismes clés}
\figslide{0.9}{figs/f5_objet}{L'héritage (une classe fille récupère et enrichit une classe mère) et le polymorphisme (un même appel donne un résultat propre à chaque type).}
\end{frame}

% #####################################################################
\section{Vous utilisez déjà des objets}

\begin{frame}{L'idée : données + comportements}
\begin{ucaddef}[L'idée en une phrase]
{\small Un \textbf{objet} réunit des \textbf{données} et les \textbf{méthodes} (fonctions attachées) qui les manipulent. Une chaîne, une liste \emph{sont} déjà des objets : les méthodes que vous appelez avec un point en sont la preuve.}
\end{ucaddef}
{\small La POO consiste à définir \textbf{nos propres} types d'objets, taillés pour notre problème — un compte bancaire, un patient, un produit — chacun avec ses données et ses comportements.}
\end{frame}

\begin{frame}[fragile]{Des objets que vous connaissez déjà}
\begin{lstlisting}
texte = "dakar"
print(texte.upper())
notes = [17, 12]
notes.append(15)
print(notes)
\end{lstlisting}
\resultat
\begin{lstlisting}[style=out]
DAKAR
[17, 12, 15]
\end{lstlisting}
{\small \texttt{.upper()} et \texttt{.append()} sont des \textbf{méthodes} : des fonctions qui appartiennent à l'objet et agissent sur lui.}
\end{frame}

% #####################################################################
\section{Définir une classe}

\begin{frame}{L'idée : la classe, plan de l'objet}
\begin{ucaddef}[L'idée en une phrase]
{\small Une \textbf{classe} est le \textbf{plan} d'un objet : elle décrit ses données et ses méthodes. \texttt{\_\_init\_\_} est le \textbf{constructeur}, exécuté à la création ; \texttt{self} y désigne l'objet en train d'être construit, sur lequel on installe les attributs avec \texttt{self.x = ...}.}
\end{ucaddef}
\end{frame}

\begin{frame}{Analogie}
\begin{ucadfront}[Le plan et l'objet]
{\small La \textbf{classe} est comme le \emph{plan d'une maison} : un seul plan, mais on peut construire autant de maisons qu'on veut. Chaque maison bâtie est une \textbf{instance} : elle suit le plan, mais possède ses propres caractéristiques (sa couleur, ses habitants). \texttt{CompteBancaire} est le plan ; \texttt{c1}, \texttt{c2} sont des comptes concrets, chacun avec son titulaire et son solde.}
\end{ucadfront}
\end{frame}

\begin{frame}[fragile]{Définir une classe et créer un objet}
\begin{lstlisting}
class CompteBancaire:
    def __init__(self, titulaire, solde):
        self.titulaire = titulaire
        self.solde = solde

c1 = CompteBancaire("Astou", 5000)
print(c1.titulaire)
print(c1.solde)
\end{lstlisting}
\resultat
\begin{lstlisting}[style=out]
Astou
5000
\end{lstlisting}
{\small \texttt{c1} est une \textbf{instance}. On accède à ses \textbf{attributs} avec un point.}
\end{frame}

\begin{frame}[fragile]{Plusieurs instances, attributs modifiables}
\begin{lstlisting}
c2 = CompteBancaire("Modou", 12000)
print(c1.titulaire, c1.solde)
print(c2.titulaire, c2.solde)
c2.solde = 20000
print(c2.solde)
\end{lstlisting}
\resultat
\begin{lstlisting}[style=out]
Astou 5000
Modou 12000
20000
\end{lstlisting}
{\small Chaque instance a ses propres valeurs ; un attribut se lit et se \textbf{modifie} avec un point.}
\end{frame}

% #####################################################################
\section{Les méthodes}

\begin{frame}{L'idée : les comportements de l'objet}
\begin{ucaddef}[L'idée en une phrase]
{\small Une \textbf{méthode} est une fonction définie \emph{dans} la classe. Elle reçoit toujours \texttt{self} en premier paramètre, ce qui lui donne accès aux attributs de l'objet (\texttt{self.solde}). Une méthode peut \textbf{modifier} l'objet, ou \textbf{renvoyer} une valeur.}
\end{ucaddef}
\end{frame}

\begin{frame}[fragile]{Des méthodes qui modifient l'objet}
\begin{lstlisting}
class CompteBancaire:
    def __init__(self, titulaire, solde):
        self.titulaire = titulaire
        self.solde = solde
    def depot(self, montant):
        self.solde = self.solde + montant
    def retrait(self, montant):
        if montant <= self.solde:
            self.solde = self.solde - montant
        else:
            print("Solde insuffisant")

c1 = CompteBancaire("Astou", 5000)
c1.depot(2000)
print(c1.solde)
c1.retrait(1000)
print(c1.solde)
\end{lstlisting}
\end{frame}

\begin{frame}[fragile]{Appel de méthode, et sortie}
{\small On appelle une méthode sur l'objet : \texttt{c1.depot(2000)}. Inutile de passer \texttt{self}, Python le fait : \texttt{self} sera \texttt{c1}.}
\resultat
\begin{lstlisting}[style=out]
7000
6000
\end{lstlisting}
{\small Un retrait supérieur au solde n'est pas effectué :}
\begin{lstlisting}
c1.retrait(50000)
print(c1.solde)
\end{lstlisting}
\begin{lstlisting}[style=out]
Solde insuffisant
6000
\end{lstlisting}
\end{frame}

\begin{frame}[fragile]{Une méthode peut renvoyer une valeur}
\begin{ucadex}[Renvoyer plutôt que modifier]
{\small \texttt{peut\_retirer} ne change rien : elle \textbf{renvoie} un booléen, utilisable dans un test.}
\end{ucadex}
\begin{lstlisting}
    def peut_retirer(self, montant):
        return montant <= self.solde
\end{lstlisting}
\begin{lstlisting}
print(c1.peut_retirer(1000))
print(c1.peut_retirer(999999))
\end{lstlisting}
\resultat
\begin{lstlisting}[style=out]
True
False
\end{lstlisting}
\end{frame}

% #####################################################################
\section{Un affichage lisible}

\begin{frame}[fragile]{La méthode spéciale \texttt{\_\_str\_\_}}
\begin{ucaddef}[L'idée en une phrase]
{\small \texttt{\_\_str\_\_} décide ce qu'affiche \texttt{print(objet)} : elle \textbf{renvoie} une chaîne. On ajoute aussi \texttt{description}, qui renvoie un libellé (utile plus loin).}
\end{ucaddef}
\begin{lstlisting}
    def __str__(self):
        return f"Compte de {self.titulaire} : {self.solde} F"
    def description(self):
        return "compte courant"
\end{lstlisting}
\begin{lstlisting}
c1 = CompteBancaire("Astou", 5000)
print(c1)
print(c1.description())
\end{lstlisting}
\resultat
\begin{lstlisting}[style=out]
Compte de Astou : 5000 F
compte courant
\end{lstlisting}
\end{frame}

% #####################################################################
\section{L'héritage}

\begin{frame}{L'idée : spécialiser sans réécrire}
\begin{ucaddef}[L'idée en une phrase]
{\small L'\textbf{héritage} crée une classe \emph{à partir d'une autre} : elle récupère tout et peut \textbf{ajouter} ce qui lui est propre. \texttt{CompteEpargne(CompteBancaire)} se lit « un compte épargne \textbf{est un} compte bancaire, spécialisé ». \texttt{super().\_\_init\_\_(...)} réutilise le constructeur du parent.}
\end{ucaddef}
\end{frame}

\begin{frame}[fragile]{Créer une classe fille}
\begin{lstlisting}
class CompteEpargne(CompteBancaire):
    def __init__(self, titulaire, solde, taux):
        super().__init__(titulaire, solde)
        self.taux = taux
    def ajouter_interets(self):
        self.solde = self.solde + self.solde * self.taux

e1 = CompteEpargne("Bineta", 10000, 0.05)
print(e1)
e1.depot(5000)
e1.ajouter_interets()
print(e1.solde)
\end{lstlisting}
\resultat
\begin{lstlisting}[style=out]
Compte de Bineta : 10000 F
15750.0
\end{lstlisting}
\end{frame}

\begin{frame}[fragile]{Tout le parent est hérité}
{\small \texttt{print(e1)} utilise le \texttt{\_\_str\_\_} \textbf{hérité} ; \texttt{e1.depot(...)} et \texttt{e1.retrait(...)} sont des méthodes \textbf{héritées}. Seuls \texttt{taux} et \texttt{ajouter\_interets} sont propres à la classe fille.}
\begin{lstlisting}
e1.retrait(2000)
print(e1.solde)
\end{lstlisting}
\resultat
\begin{lstlisting}[style=out]
13750.0
\end{lstlisting}
{\small Le retrait, écrit une seule fois dans \texttt{CompteBancaire}, fonctionne sur le compte épargne sans une ligne de plus.}
\end{frame}

% #####################################################################
\section{Le polymorphisme}

\begin{frame}{L'idée : un même appel, des résultats différents}
\begin{ucaddef}[L'idée en une phrase]
{\small \textbf{Surcharger} une méthode, c'est la \textbf{redéfinir} dans la classe fille. Le \textbf{polymorphisme} en découle : un \emph{même} appel de méthode produit un résultat \textbf{propre au type réel} de l'objet. Python choisit automatiquement la bonne version.}
\end{ucaddef}
\end{frame}

\begin{frame}[fragile]{Surcharger une méthode}
\begin{lstlisting}
class CompteEpargne(CompteBancaire):
    def __init__(self, titulaire, solde, taux):
        super().__init__(titulaire, solde)
        self.taux = taux
    def description(self):
        return f"compte epargne a {self.taux}"

comptes = [CompteBancaire("Astou", 5000), CompteEpargne("Bineta", 10000, 0.05)]
for c in comptes:
    print(c.titulaire, ":", c.description())
\end{lstlisting}
\resultat
\begin{lstlisting}[style=out]
Astou : compte courant
Bineta : compte epargne a 0.05
\end{lstlisting}
\end{frame}

\begin{frame}[fragile]{Le polymorphisme à travers une fonction}
\begin{lstlisting}
def presenter(compte):
    print("Il s'agit d'un", compte.description())

presenter(comptes[0])
presenter(comptes[1])
\end{lstlisting}
\resultat
\begin{lstlisting}[style=out]
Il s'agit d'un compte courant
Il s'agit d'un compte epargne a 0.05
\end{lstlisting}
{\small La fonction \texttt{presenter} accepte n'importe quel compte : elle appelle \texttt{description()} sans se soucier du type exact. C'est toute la souplesse du polymorphisme.}
\end{frame}

% #####################################################################
\section{Synthèse}

\begin{frame}{Pièges fréquents}
\begin{itemize}\small
  \item \textbf{Oublier \texttt{self}} : toute méthode le prend en premier paramètre, et l'on écrit \texttt{self.attribut} pour lire les données de l'objet.
  \item \textbf{Confondre la classe et l'instance} : \texttt{CompteBancaire} est le plan ; \texttt{c1 = CompteBancaire(...)} est un objet concret. On appelle les méthodes sur l'\textbf{instance}.
  \item \textbf{Méthode sans parenthèses} : \texttt{c1.description} seul désigne la méthode ; il faut \texttt{c1.description()} pour l'\textbf{exécuter}.
  \item \textbf{Oublier \texttt{super().\_\_init\_\_}} dans une classe fille : les attributs du parent ne sont alors pas installés.
\end{itemize}
\end{frame}

\begin{frame}{À retenir — Séance 5}
\begin{ucadret}[L'essentiel]
{\small
\textbf{Classe} : le plan d'un objet (\texttt{class}). \textbf{Instance} : un objet concret.\\[2pt]
\textbf{\_\_init\_\_(self, ...)} : le constructeur ; installe les \textbf{attributs} (\texttt{self.x = ...}).\\[2pt]
\textbf{self} : l'objet courant ; toute méthode le reçoit en premier.\\[2pt]
\textbf{Méthode} : fonction dans la classe, appelée par \texttt{objet.methode(...)} ; modifie l'objet ou \textbf{renvoie} une valeur. \textbf{\_\_str\_\_} : la chaîne affichée par \texttt{print}.\\[2pt]
\textbf{Héritage} : \texttt{class Fille(Parent)} récupère tout et ajoute ; \texttt{super().\_\_init\_\_(...)} réutilise le constructeur parent.\\[2pt]
\textbf{Polymorphisme} : une méthode \textbf{surchargée} donne, pour un même appel, un comportement propre à chaque classe.}
\end{ucadret}
\end{frame}

\begin{frame}{Prochaine séance}
\begin{ucaddef}[Séance 6 — NumPy]
{\small On sait organiser données et comportements en objets. La Séance 6 ouvre la \textbf{science des données} avec \texttt{numpy} : le \textbf{tableau} (\texttt{array}), le calcul \textbf{vectorisé} (une opération sur tout le tableau, sans boucle) et les agrégations. C'est l'outil sur lequel repose \texttt{pandas}.}
\end{ucaddef}
{\small À faire d'ici là : refaire les cellules « À vous » du notebook de la séance.}
\end{frame}

\begin{frame}{Références}
\begin{itemize}\small
  \item Severance, C. — \emph{Python for Everybody}, 2016.
  \item Downey, A. — \emph{Think Python}, 2e éd., O'Reilly, 2015.
  \item Matthes, E. — \emph{Python Crash Course}, 3e éd., No Starch Press, 2023.
  \item Python Software Foundation — \emph{The Python Tutorial} (Classes), 2024.
\end{itemize}
\end{frame}

\end{document}
